\documentclass[12pt]{article}
\usepackage[utf8]{inputenc}
\usepackage{amsmath,amssymb,amsthm}
\usepackage{hyperref}
\usepackage{graphicx}
\usepackage{booktabs}
\usepackage[margin=1in]{geometry}

\newtheorem{theorem}{Theorem}
\newtheorem{definition}[theorem]{Definition}

\title{First Formally Verified Braided Modular Tensor Categories\\and Knot Invariants in Lean~4}

\author{John Roehm}
\date{Draft --- April 2026}

\begin{document}

\maketitle

\begin{abstract}
We present the first machine-checked formalization of braided modular tensor
categories (MTCs), categorical knot invariants, and topological quantum field
theory (TQFT) partition functions in any proof assistant. Working in Lean~4
with Mathlib, we construct verified instances of the Ising and Fibonacci MTCs
--- including complete F-symbol data, pentagon and hexagon equations, ribbon
conditions, and S-matrix unitarity --- all proved by \texttt{native\_decide}
over custom decidable algebraic number fields. As applications, we compute
the first verified knot invariants from categorical data: the trefoil
evaluates to $-1$ (Jones polynomial at $q = e^{2\pi i/16}$) and the
figure-eight knot to $1 - q^2 - q^{-2}$. We further verify TQFT partition
functions on surfaces of genus $g = 0, 1, 2$ for both Ising and Fibonacci
theories. The formalization comprises % Auto-generated by update_counts.py — 2026-04-24T19:44:03.140400
% Do not edit manually. Run: uv run python scripts/update_counts.py
\newcommand{\totaltheorems}{3993}
\newcommand{\substantivetheorems}{3882}
\newcommand{\placeholdertheorems}{111}
\newcommand{\axiomcount}{1}
\newcommand{\sorrycount}{0}
\newcommand{\sorrypercent}{0.0}
\newcommand{\leanmodules}{167}
\newcommand{\leandefinitions}{3297}
\newcommand{\pythonmodules}{68}
\newcommand{\testfiles}{59}
\newcommand{\totaltests}{2836}
\newcommand{\figurecount}{111}
\newcommand{\notebookcount}{54}
\newcommand{\papercount}{19}
\newcommand{\aristotleproved}{322}
\newcommand{\aristotleruns}{44}
% --- Paper 18 / Phase 5t per-section counts (FermiHubbardDimer.lean) ---
\newcommand{\fhdTotal}{143}
\newcommand{\fhdWaveTwo}{13}
\newcommand{\fhdWaveThree}{6}
\newcommand{\fhdWaveFour}{22}
\newcommand{\fhdWaveFive}{22}
\newcommand{\fhdWaveSixABC}{38}
\newcommand{\fhdWaveSixExtra}{15}
\newcommand{\fhdWaveSeven}{19}
\newcommand{\fhdWaveEight}{8}
\totaltheorems{}
theorems across % Auto-generated by update_counts.py — 2026-04-24T19:44:03.140400
% Do not edit manually. Run: uv run python scripts/update_counts.py
\newcommand{\totaltheorems}{3993}
\newcommand{\substantivetheorems}{3882}
\newcommand{\placeholdertheorems}{111}
\newcommand{\axiomcount}{1}
\newcommand{\sorrycount}{0}
\newcommand{\sorrypercent}{0.0}
\newcommand{\leanmodules}{167}
\newcommand{\leandefinitions}{3297}
\newcommand{\pythonmodules}{68}
\newcommand{\testfiles}{59}
\newcommand{\totaltests}{2836}
\newcommand{\figurecount}{111}
\newcommand{\notebookcount}{54}
\newcommand{\papercount}{19}
\newcommand{\aristotleproved}{322}
\newcommand{\aristotleruns}{44}
% --- Paper 18 / Phase 5t per-section counts (FermiHubbardDimer.lean) ---
\newcommand{\fhdTotal}{143}
\newcommand{\fhdWaveTwo}{13}
\newcommand{\fhdWaveThree}{6}
\newcommand{\fhdWaveFour}{22}
\newcommand{\fhdWaveFive}{22}
\newcommand{\fhdWaveSixABC}{38}
\newcommand{\fhdWaveSixExtra}{15}
\newcommand{\fhdWaveSeven}{19}
\newcommand{\fhdWaveEight}{8}
\totalmodules{} Lean modules with
zero \texttt{sorry} in the MTC/knot/TQFT components. The decidable number
field technique --- representing elements of cyclotomic fields as rational
coefficient tuples with computable arithmetic --- appears to be novel and
applicable to algebraic verification tasks beyond MTCs.
\end{abstract}

% ============================================================
\section{Introduction}
\label{sec:intro}

% Frame: first-ever, three audiences, why it matters
% - Formalization community: novel Lean 4 definitions (PivotalCategory, RibbonCategory, FusionCategory, MTC)
% - Mathematical physics: certified algebraic data for topological quantum computing
% - TQC community: verified anyon braiding foundations (Majorana 1 controversy context)

% Key claims:
% 1. First PivotalCategory, RibbonCategory, FusionCategory, MTC in any proof assistant
% 2. First pentagon/hexagon equations verified for specific anyon models
% 3. First categorical knot invariants computed in a proof assistant
% 4. Novel decidable number field technique for algebraic verification
% 5. Zero sorry across all MTC/knot/TQFT modules

% ============================================================
\section{Background: Modular Tensor Categories}
\label{sec:background}

% Brief mathematical background
% - Fusion categories, F-symbols, pentagon equation
% - Braided categories, R-matrix, hexagon equations
% - Ribbon categories, twist, Gauss sum
% - Modular tensor categories, S-matrix, non-degeneracy

% ============================================================
\section{Lean~4 Formalization Architecture}
\label{sec:architecture}

% Typeclass hierarchy:
%   MonoidalCategory (Mathlib) -> BraidedCategory (Mathlib) -> RigidCategory (Mathlib)
%   -> PivotalCategory (ours) -> SphericalCategory (ours) -> RibbonCategory (ours)
%   -> PreModularData (ours) -> ModularTensorData (ours, non-degenerate S-matrix)
%
% FusionCategoryData: object type, fusion coefficients, F-symbols, pentagon
% BraidedFusionData: adds R-matrix, hexagon equations
%
% Concrete instances via decidable types + native_decide

% ============================================================
\section{Decidable Number Fields for Algebraic Verification}
\label{sec:number_fields}

% The technique:
% - Q(sqrt(2)) for Ising F-symbols (degree 2)
% - Q(zeta_5) for Fibonacci (degree 4)
% - Q(zeta_16) for Ising R-matrix (degree 8)
% - Q(sqrt(3)), Q(sqrt(5)) for higher-k S-matrices
%
% Pattern: Q[x]/(f(x)) as rational tuple with computable reduction
% native_decide: compiler-trusted computation over DecidableEq types
% Trust trade-off: ~30K lines trusted (compiler) vs kernel-only (norm_num)

% ============================================================
\section{Verified MTC Instances}
\label{sec:instances}

\subsection{Ising MTC (SU(2)$_2$)}
% F-symbols: 36 nonzero entries in Kitaev convention
% Pentagon: all 3^5 = 243 cases verified by native_decide over Q(sqrt(2))
% R-matrix: 6 nonzero entries in Q(zeta_16)
% Hexagon: all 6 equations verified
% Ribbon: 4 conditions verified
% S-matrix: 3x3, unitarity S*S^T = I verified
% Gauss sum: tau_+ * tau_- = D^2

\subsection{Fibonacci MTC}
% F-symbols: golden ratio phi in Q(sqrt(5))
% Pentagon: 2^5 = 32 relevant cases (2 objects) verified
% R-matrix: 2 entries in Q(zeta_5)
% Hexagon: 3 equations verified
% Trefoil knot invariant: -1 (FIRST verified knot invariant)

\subsection{SU(2)$_k$ at Higher Levels}
% k=3: 4 objects, S-matrix in Q(level3)
% k=4: 5 objects, S-matrix in Q(sqrt(3))
% Unitarity verified for all levels k=1,...,4

% ============================================================
\section{Knot Invariants from Categorical Data}
\label{sec:knots}

% BraidGroup via PresentedGroup + relations
% RT representation: braid -> endomorphism in MTC
% Trefoil (sigma_1^3): trace = -1 in Ising (PROVED)
% Hopf link (sigma_1^2): trace = 0 in Fibonacci (PROVED)
% Figure-eight (sigma_1 sigma_2^{-1} sigma_1 sigma_2^{-1}): computed

% ============================================================
\section{TQFT Partition Functions}
\label{sec:tqft}

% Z(S^2) = 1, Z(T^2) = n (number of anyons)
% Z(Sigma_g) = sum_a d_a^{2-2g} (Verlinde formula)
% Verified for Ising (g=0,1,2) and Fibonacci (g=0,1,2)

% ============================================================
\section{SU(3)$_k$ Fusion Categories}
\label{sec:su3}

% NEW: first SU(3)_k formalization
% SU(3)_1: Z_3 fusion ring, 3 objects
% SU(3)_2: 6 anyons, Fibonacci subcategory (tau x tau = 1 + tau)
% All fusion rules verified by native_decide
% Associativity and commutativity proved

% ============================================================
\section{Discussion}
\label{sec:discussion}

% What the formalization revealed
% - Kitaev convention for F-symbols is the natural one
% - Pentagon convention mismatch was a real bug (caught by verification)
% - native_decide over custom fields is dramatically more practical than norm_num
% - Zero sorry achieved across all MTC/knot/TQFT modules

% Limitations:
% - native_decide trusts the compiler
% - Custom number fields are hand-written (consolidation in progress)
% - No pentagon for SU(3)_k yet (requires F-symbols in Q(zeta_15))

% ============================================================
\section{Related Work}
\label{sec:related}

% Schreiber/Sati/Myers: HoTT topological quantum gates (Comm. Math. Phys. 2024)
% SQIR/VOQC: verified quantum circuits (POPL 2021)
% Lean-QuantumInfo: quantum Stein's lemma (Perimeter Institute)
% PhysLean (HepLean): quantum harmonic oscillator, Higgs potential
% Mathlib: braided monoidal, rigid categories (abstract, not concrete instances)

% ============================================================
\section{Conclusion}
\label{sec:conclusion}

% Summary of firsts
% Future: Mathlib PR for categorical infrastructure, more MTC instances
% The decidable number field technique is general-purpose

\bibliographystyle{plain}
\bibliography{refs}

\end{document}
